\documentclass[12pt,a4paper]{article}

\usepackage[utf8]{inputenc}
\usepackage[margin=1in]{geometry}
\usepackage{amsmath,amssymb}
\usepackage{graphicx}
\usepackage{booktabs}
\usepackage{listings}
\usepackage{xcolor}
\usepackage{hyperref}
\usepackage{microtype}
\usepackage{caption}
\usepackage{setspace}
\usepackage{tikz}
\usetikzlibrary{arrows.meta,positioning,shapes.geometric,fit,backgrounds}

\hypersetup{
    colorlinks=true,
    linkcolor=green!40!black,
    citecolor=blue!70!black,
    urlcolor=teal!70!black,
    pdftitle={AI-Powered Eco-Habit and Green Footprint Coach},
    pdfauthor={Prince Jaiswal}
}

% Code listing styling matching modern IDE aesthetic
\definecolor{codebg}{rgb}{0.97,0.98,0.98}
\definecolor{codeborder}{rgb}{0.85,0.88,0.85}
\definecolor{keywordcolor}{rgb}{0.08,0.40,0.20}
\definecolor{stringcolor}{rgb}{0.12,0.50,0.30}
\definecolor{commentcolor}{rgb}{0.45,0.50,0.45}

\lstdefinelanguage{JavaScript}{
    keywords={const,let,var,await,async,function,return,if,else,for,of,new,import,from,export,try,catch,switch,case,break,default},
    sensitive=true,
    comment=[l]{//},
    string=[b]"
}
\lstset{
    backgroundcolor=\color{codebg},
    frame=single,
    rulecolor=\color{codeborder},
    basicstyle=\ttfamily\footnotesize,
    breaklines=true,
    captionpos=b,
    tabsize=2,
    showstringspaces=false,
    keywordstyle=\color{keywordcolor}\bfseries,
    stringstyle=\color{stringcolor},
    commentstyle=\color{commentcolor}\itshape
}

\onehalfspacing

\begin{document}

% =========================================================================
% FRONT MATTER: title, author, abstract
% =========================================================================
\title{\textbf{AI-Powered Eco-Habit \& Green Footprint Coach:\\A Full-Stack Sustainability Platform with Conversational NLP\\and Gamified Behavioral Intervention}}

\author{%
\begin{minipage}[t]{0.46\textwidth}\centering
\textbf{Prince Jaiswal$^1$}\\
\small Department of Computer Applications\\
\small BCA Final Year Project
\end{minipage}
\hfill
\begin{minipage}[t]{0.46\textwidth}\centering
\textbf{Dr. Ravindra Nath$^2$}\\
\small Project Supervisor / Guide\\
\small Faculty of Computer Applications
\end{minipage}}
\date{}
\maketitle

\begin{abstract}
Mitigating individual greenhouse gas (GHG) emissions is a vital pillar in achieving global net-zero targets established by the Paris Climate Accord. However, conventional personal carbon accounting platforms present significant adoption barriers: users are confronted with tedious, multi-page questionnaires, complex technical units, static one-off calculations, and generic advice devoid of continuous coaching. This academic project presents the design and implementation of the \textbf{AI-Powered Eco-Habit \& Green Footprint Coach}, a production-ready, full-stack web platform combining conversational Natural Language Processing (NLP), internationally recognized carbon accounting methodologies, dynamic data analytics, and behavioral gamification.

The platform's primary contribution is an \textbf{AI-Powered Natural Language Activity Logger} leveraging the Google Gemini API. Users input arbitrary, natural descriptions of daily routines (e.g., \emph{``Today I travelled 15 km by car, used AC for 5 hours, ate chicken for lunch, and used public transport for 10 km''}). The language model interprets the input, extracts discrete activity items, applies scientifically documented emission factors from DEFRA (UK), EPA (US), and the IPCC, and returns structured, schema-validated JSON with explicit assumptions and follow-up prompts for missing parameters. The architecture incorporates an offline deterministic NLP engine ensuring zero downtime if external APIs are unconfigured. 

Complementing conversational logging, the application features an interactive multi-category carbon calculator, daily aggregated telemetry indexing (\texttt{CarbonRecord}) powering responsive Recharts visualizations across 7-, 30-, and 90-day intervals, an adaptive AI recommendation engine analyzing 30-day historical footprint patterns, and an Octalysis-inspired gamification system incorporating Eco Points, progressive user tiers (\emph{Seedling} to \emph{Earth Champion}), daily logging streaks, milestone badges, and community challenges. Built on a decoupled React.js (Vite), Express.js, and MongoDB architecture with JWT authentication, the platform demonstrates that modern web frameworks and generative AI can eliminate cognitive friction in personal carbon accounting and foster long-term sustainable behavior.
\end{abstract}

\noindent\textbf{Keywords:} Personal Carbon Accounting, Conversational AI, Natural Language Processing, Large Language Models, Google Gemini API, DEFRA Conversion Factors, MERN Stack, Recharts, Gamification, Behavioral Sustainability.

% =========================================================================
% 1. INTRODUCTION
% =========================================================================
\section{Introduction}
Anthropogenic climate change represents one of the most critical challenges of the modern era. Climate science dictates that limiting global temperature increases to within $1.5^\circ\text{C}$ requires atmospheric greenhouse gas emissions to drop precipitously by 2030, culminating in global carbon neutrality by mid-century. While international treaties and industrial regulations govern macro-level emissions, individual household consumption directly and indirectly accounts for over 65\% of global emissions through transportation choices, residential energy usage, dietary composition, and consumer goods consumption.

To reach the $1.5^\circ\text{C}$ trajectory, per-capita emissions must converge toward a sustainable baseline of approximately $5.0\text{ kg CO}_2\text{e per day}$ ($\approx 1.8$ to $2.0\text{ tonnes per year}$). In contrast, current average daily footprints in industrialized and emerging nations frequently exceed $25\text{ to }40\text{ kg CO}_2\text{e}$. Bridging this gap requires empowering individuals with real-time awareness and actionable insights into their everyday environmental impact.

\subsection{Problem Statement}
Existing carbon calculator solutions suffer from three fundamental limitations:
\begin{enumerate}
    \item \textbf{Cognitive Friction and Manual Data Entry:} Most platforms mandate extensive questionnaires requiring exact kilowatt-hours, vehicle fuel economy, and numerical weights. This causes severe user fatigue and high drop-off rates.
    \item \textbf{Static, Non-Continuous Accounting:} Users calculate their footprint once as a novelty rather than establishing continuous daily logging and tracking.
    \item \textbf{Lack of Actionable, Contextual Coaching:} Traditional outputs produce abstract totals (e.g., ``You emit 12 tonnes/year'') without providing granular habit interventions, personalized reduction trajectories, or motivational incentives.
\end{enumerate}

\subsection{Objectives}
The primary objective of this project is to architect, implement, test, and evaluate a comprehensive full-stack platform named \textbf{AI-Powered Eco-Habit \& Green Footprint Coach}. The specific technical objectives are:
\begin{enumerate}
    \item \textbf{Conversational NLP Logging:} Implement a zero-friction logging mechanism where natural English prompts are parsed by Google Gemini AI into discrete, categorized activities with estimated $\text{kg CO}_2\text{e}$ emissions.
    \item \textbf{Scientific Emission Estimation:} Base calculations on published, internationally recognized greenhouse gas conversion factors (DEFRA, EPA, and IPCC) rather than arbitrary metrics, with explicit accounting for assumptions and passenger sharing.
    \item \textbf{High-Performance Telemetry Visualizations:} Maintain normalized daily records in MongoDB to feed real-time Recharts visualizations (area trends, weekly bars, and category pie charts) over 7, 30, and 90-day intervals.
    \item \textbf{Personalized Green Habit Recommendations:} Engineer an AI recommendation pipeline that diagnoses a user's 30-day telemetry to surface high-impact, achievable lifestyle changes with estimated carbon savings.
    \item \textbf{Gamified Behavioral Motivation:} Implement an anti-abuse gamification layer featuring Eco Points, five user levels, consecutive streaks, achievement badges, and weekly eco challenges with a community leaderboard.
    \item \textbf{Robust Zero-Docker Deployment:} Provide seamless execution on local and cloud environments through in-memory MongoDB fallback and deterministic NLP fallback engines.
\end{enumerate}

\subsection{Scope of the Project}
This system is developed as an academic BCA final-year project. It covers individual lifestyle carbon accounting across six core categories: transportation, residential energy and air conditioning, food and diet, consumer shopping, domestic water usage, and municipal waste management. The platform features responsive user interfaces tailored for mobile, tablet, and desktop viewports, complete with secure JWT user authentication and private multi-tenant data isolation.

% =========================================================================
% 2. LITERATURE REVIEW
% =========================================================================
\section{Literature Review}
\textbf{Personal Carbon Accounting (PCA).} Research into personal carbon accounting frameworks highlights the necessity of real-time feedback loops. Early systems relied on static life-cycle assessment (LCA) spreadsheets. Studies by West et al.\ and Lin et al.\ demonstrated that automated, low-latency footprint tracking produces up to a 14\% sustained reduction in domestic energy consumption compared to quarterly utility bill reviews.

\textbf{Natural Language Processing for Environmental Sensing.} Recent breakthroughs in Large Language Models (LLMs) demonstrate state-of-the-art capability in extracting structured tabular entities from unstructured prose. Rather than forcing users through rigid form inputs, schema-guided prompt engineering enables models like Google Gemini to identify quantities, modalities, durations, and dietary components directly from conversational natural language.

\textbf{Documented Greenhouse Gas Conversion Standards.} Accurate personal carbon accounting necessitates rigorous adherence to recognized international standards:
\begin{itemize}
    \item \emph{UK Department for Environment, Food \& Rural Affairs (DEFRA):} Publishes annual government GHG conversion factors for company reporting, providing standardized per-passenger-km factors for passenger cars ($0.170\text{ kg CO}_2\text{e/km}$ petrol, $0.165\text{ kg CO}_2\text{e/km}$ diesel), public transit ($0.089\text{ kg CO}_2\text{e/km}$ bus, $0.035\text{ kg CO}_2\text{e/km}$ metro rail), and domestic aviation ($0.246\text{ kg CO}_2\text{e/km}$).
    \item \emph{US Environmental Protection Agency (EPA):} Provides greenhouse gas equivalency methodologies and grid electricity emission baselines ($\approx 0.475\text{ kg CO}_2\text{e/kWh}$ average).
    \item \emph{Poore \& Nemecek / IPCC Guidelines:} Landmark agricultural meta-analyses detailing lifecycle dietary carbon intensities (e.g., beef at $27.0\text{ kg CO}_2\text{e/kg}$ food vs.\ poultry at $5.2\text{ kg CO}_2\text{e/kg}$ and legumes at $0.9\text{ kg CO}_2\text{e/kg}$).
\end{itemize}

\textbf{Nudge Theory and Gamification in Pro-Environmental Behavior.} Thaler and Sunstein's nudge theory emphasizes structuring choice architectures to encourage positive decisions without mandates. Incorporating gamification elements—such as points, badges, progress bars, and social comparisons—leverages intrinsic and extrinsic motivation to reinforce daily habits.

% =========================================================================
% 3. SYSTEM ARCHITECTURE
% =========================================================================
\section{System Architecture}
The system is architected as a three-tier decoupled web application comprising a single-page React client, an Express.js REST API gateway, external AI/database services, and an embedded fallback engine (Figure~\ref{fig:arch}).

\begin{figure}[htbp]
\centering
\begin{tikzpicture}[
    node distance=1.1cm and 1.5cm,
    box/.style={draw, rounded corners, align=center, minimum width=3.4cm, minimum height=1.0cm, fill=emerald!5},
    db/.style={draw, cylinder, shape border rotate=90, aspect=0.25, align=center, minimum width=3.4cm, minimum height=1.3cm, fill=green!10},
    ext/.style={draw, rounded corners, align=center, minimum width=3.2cm, minimum height=1.0cm, fill=amber!10},
    arr/.style={-{Stealth[length=2.5mm]}, thick}
]
\node[box] (fe) {React 18 + Vite SPA\\(Port 5173)\\\small Tailwind + Recharts};
\node[box, right=of fe] (gw) {Express.js API Gateway\\(Port 5001)\\\small Auth, Routes, Controllers};
\node[box, below=1.4cm of gw] (calc) {Calculation \& NLP\\Services\\\small DEFRA / EPA / IPCC};
\node[db, below=1.4cm of fe] (mongo) {MongoDB Database\\(Atlas / In-Memory)\\\small 7 Mongoose Models};
\node[ext, right=of gw] (gemini) {Google Gemini AI\\(\texttt{gemini-1.5-flash})\\\small Structured JSON Schema};

\draw[arr] (fe) -- node[above, font=\scriptsize]{REST / JWT} (gw);
\draw[arr] (gw) -- node[right, font=\scriptsize]{Internal Dispatch} (calc);
\draw[arr] (gw) -- node[above, font=\scriptsize]{Prompt} (gemini);
\draw[arr] (gemini) -- node[below, font=\scriptsize]{Structured JSON} (gw);
\draw[arr] (gw) -- node[left, font=\scriptsize]{Mongoose ODM} (mongo);
\draw[arr] (calc) -- node[below, font=\scriptsize]{Sync Records} (mongo);
\end{tikzpicture}
\caption{System architecture of the Eco-Habit \& Green Footprint Coach.}
\label{fig:arch}
\end{figure}

The platform comprises the following decoupled modules:
\begin{itemize}
    \item \textbf{Client Tier (React 18, Vite 6, Tailwind CSS):} Provides a modern SaaS dashboard inspired by Linear and Stripe. Manages authentication sessions, handles optimistic UI updates, renders responsive Recharts diagrams, and hosts interactive logging interfaces.
    \item \textbf{Backend Gateway (Express.js 4, Node.js):} Implements centralized routing, JWT verification middleware, input validation, Helmet security headers, and rate limiting (300 requests/15 min for general API, 60 requests/5 min for AI endpoints).
    \item \textbf{Scientific Calculation Service:} Houses standardized DEFRA, EPA, and IPCC mathematical models, computing transport, energy, food, shopping, water, and waste emissions deterministically.
    \item \textbf{AI Service (Google Gemini API):} Analyzes natural language strings using a zero-shot environmental accounting system instruction, returning verified JSON structures with discrete activity decompositions.
    \item \textbf{Persistence Layer (MongoDB + Mongoose 8):} Stores users, activities, normalized daily telemetry records, recommendations, challenges, user progress, and badges with indexing on \texttt{\{userId, dateString\}}.
\end{itemize}

% =========================================================================
% 4. METHODOLOGY
% =========================================================================
\section{Methodology}

\subsection{Natural Language Processing and Extraction Pipeline}
When a user submits a freeform activity log, the backend pipeline executes nine sequential stages:
\begin{enumerate}
    \item \textbf{Input Sanitization:} Trims whitespace, validates minimum prompt length, and verifies JWT identity.
    \item \textbf{Prompt Construction:} Combines the user's natural language input with a strict system prompt containing DEFRA/EPA guidelines and a target JSON schema.
    \item \textbf{Generative Inference:} Dispatches the payload to the Gemini API (`gemini-1.5-flash`) with \texttt{responseMimeType: "application/json"} and temperature $0.2$ to ensure deterministic outputs.
    \item \textbf{Fallback NLP Handler:} If \texttt{GEMINI\_API\_KEY} is unconfigured or encounters a network outage, a regex-based deterministic extraction engine automatically parses known quantities (e.g., kilometers, hours, meal types) and assigns standardized factors.
    \item \textbf{Schema Validation:} Strips potential markdown fences, parses the JSON payload, verifies non-empty activity arrays, and enforces numeric constraints ($E_i \ge 0$).
    \item \textbf{Assumption Verification:} Validates that assumed quantities (such as car fuel types or AC power) are explicitly annotated in an \texttt{assumptions} field.
    \item \textbf{Database Persistence:} Writes an \texttt{Activity} record linking discrete items, summary text, and eco scores to the authenticated user.
    \item \textbf{Telemetry Aggregation:} Upserts the daily \texttt{CarbonRecord} for the specified date string (\texttt{YYYY-MM-DD}), synchronizing category totals.
    \item \textbf{Gamification Processing:} Updates consecutive daily streak counters, awards $+15$ Eco Points, and evaluates milestone achievement rules.
\end{enumerate}

\subsection{Mathematical Emission Formulations}
Emissions are calculated in kilograms of carbon dioxide equivalent ($\text{kg CO}_2\text{e}$) across six distinct lifestyle domains:

\paragraph{1. Transportation.}
For distance $d$ (km), passenger vehicle mode factor $f_{\text{mode}}$, and passenger count $p \ge 1$:
\begin{equation}
    E_{\text{trans}} = \frac{d \cdot f_{\text{mode}}}{p}
\end{equation}
where $f_{\text{petrol}} = 0.170$, $f_{\text{diesel}} = 0.165$, $f_{\text{hybrid}} = 0.105$, $f_{\text{ev}} = 0.047$, $f_{\text{bus}} = 0.089$, $f_{\text{metro}} = 0.035$, and $f_{\text{bike}} = f_{\text{walk}} = 0.000\text{ kg CO}_2\text{e/km}$.

\paragraph{2. Electricity and HVAC.}
For operating hours $t$ and appliance rated hourly consumption factor $f_{\text{app}}$:
\begin{equation}
    E_{\text{energy}} = t \cdot f_{\text{app}}
\end{equation}
For direct metered consumption in kilowatt-hours ($k$):
\begin{equation}
    E_{\text{energy}} = k \cdot I_{\text{grid}} \quad (I_{\text{grid}} = 0.475\text{ kg CO}_2\text{e/kWh})
\end{equation}
where standard appliance coefficients include $f_{\text{AC}} = 0.50\text{ kg/hr}$ (1.5 kW compressor), $f_{\text{heater}} = 0.95\text{ kg/hr}$ (2.0 kW element), $f_{\text{desktop}} = 0.07\text{ kg/hr}$ (150 W), $f_{\text{laptop}} = 0.02\text{ kg/hr}$ (40 W), and $f_{\text{geyser}} = 1.50\text{ kg/hr}$ (3.0 kW).

\paragraph{3. Food and Dietary Consumption.}
For servings $s$ and dietary classification coefficient $f_{\text{meal}}$:
\begin{equation}
    E_{\text{food}} = s \cdot f_{\text{meal}}
\end{equation}
with $f_{\text{beef}} = 6.61\text{ kg}$, $f_{\text{chicken}} = 1.30\text{ kg}$, $f_{\text{fish}} = 1.10\text{ kg}$, $f_{\text{vegetarian}} = 0.85\text{ kg}$, and $f_{\text{vegan}} = 0.55\text{ kg CO}_2\text{e/serving}$.

\paragraph{4. Water Usage.}
For shower duration $m$ (minutes) or volume $v$ (liters):
\begin{equation}
    E_{\text{water}} = \frac{m}{10} \cdot 0.45 \quad \text{or} \quad E_{\text{water}} = v \cdot 0.0003
\end{equation}

\paragraph{5. Waste and Consumer Shopping.}
For landfill waste weight $w$ (kg) and purchased goods count $n$:
\begin{align}
    E_{\text{waste}} &= w \cdot 0.58\text{ kg/kg (landfill)} \quad \text{vs.} \quad w \cdot 0.08\text{ kg/kg (recycled)} \\
    E_{\text{shopping}} &= n \cdot f_{\text{item}} \quad (f_{\text{shirt}} = 4.50, f_{\text{jeans}} = 20.00, f_{\text{plastic\_bottle}} = 0.082\text{ kg})
\end{align}

\paragraph{Total Daily Carbon Footprint.}
The total daily footprint represents the direct sum of all discrete parsed activities:
\begin{equation}
    E_{\text{total}} = \sum_{i=1}^{M} E_i = E_{\text{trans}} + E_{\text{energy}} + E_{\text{food}} + E_{\text{shopping}} + E_{\text{water}} + E_{\text{waste}}
\end{equation}

\subsection{Normalized Eco Score Algorithm}
To translate raw kilograms into an intuitive consumer benchmark, the platform computes a non-linear Eco Score $S_{\text{eco}} \in [0, 100]$:
\begin{equation}
    S_{\text{eco}} = 
    \begin{cases}
        98, & E_{\text{total}} \le 0.1 \\
        92, & 0.1 < E_{\text{total}} \le 3.0 \\
        82, & 3.0 < E_{\text{total}} \le 6.0 \\
        68, & 6.0 < E_{\text{total}} \le 10.0 \\
        52, & 10.0 < E_{\text{total}} \le 15.0 \\
        \max\left(12, 100 - 2.5 \cdot E_{\text{total}}\right), & E_{\text{total}} > 15.0
    \end{cases}
\end{equation}
Scores $\ge 80$ indicate high sustainability aligned with Paris goals, scores between $50\text{--}79$ represent moderate footprints, and scores $<50$ signal severe carbon-intensive habits requiring remediation.

\subsection{Daily Telemetry Aggregation and Compound Indexing}
Querying raw activity histories for 90-day time-series plots introduces unacceptable aggregation latency on un-indexed collections. The system solves this via the \texttt{CarbonRecord} collection, which maintains daily totals synchronized via the \texttt{syncDailyCarbonRecord} routine. A unique compound index on \texttt{\{userId: 1, dateString: 1\}} guarantees idempotency and sub-millisecond retrieval.

\subsection{Gamification and Fraud-Proof Reward Protocols}
To foster enduring behavior change, the platform implements four gamified reward mechanisms:
\begin{enumerate}
    \item \textbf{User Progression Tiers:} Five progressive levels mapped to cumulative Eco Points: Seedling ($0\text{--}99$), Sprout ($100\text{--}299$), Eco Scout ($300\text{--}599$), Green Guardian ($600\text{--}999$), and Earth Champion ($1{,}000+$).
    \item \textbf{Consecutive Day Streaks:} Evaluated by calculating calendar day differences between the current transaction and \texttt{lastActiveDate}. Day $+1$ increments the streak; day $>1$ resets the counter to 1.
    \item \textbf{Duplicate-Proof Challenge Completion:} When a user advances a challenge to its target duration, the backend checks the \texttt{pointsAwarded} boolean flag. Reward points are disbursed exactly once, preventing replay attacks.
    \item \textbf{Dynamic Habit Adoption:} Adopting an AI green recommendation marks its status as \texttt{completed} and awards $+25$ points atomically.
\end{enumerate}

% =========================================================================
% 5. INPUTS, PROCESSING AND OUTPUTS
% =========================================================================
\section{Inputs, Processing and Outputs}

\subsection{System Inputs and Endpoints}
The platform exposes RESTful endpoints supporting both conversational AI and manual calculator workflows (Table~\ref{tab:endpoints}).

\begin{table}[htbp]
\centering
\small
\begin{tabular}{@{}p{1.8cm}p{4.2cm}p{3.2cm}p{4.8cm}@{}}
\toprule
\textbf{Method} & \textbf{Endpoint} & \textbf{Auth Required} & \textbf{Function} \\ \midrule
\texttt{POST} & \texttt{/api/auth/register} & No & Registers user; awards $+50$ welcome pts \\ \addlinespace
\texttt{POST} & \texttt{/api/auth/login} & No & Authenticates user; returns JWT token \\ \addlinespace
\texttt{POST} & \texttt{/api/activities/analyze} & Yes (JWT) & Parses NL prompt using Gemini AI \\ \addlinespace
\texttt{POST} & \texttt{/api/activities/manual} & Yes (JWT) & Saves multi-category calculator tray \\ \addlinespace
\texttt{GET}  & \texttt{/api/analytics/summary} & Yes (JWT) & Fetches top stat cards and change \% \\ \addlinespace
\texttt{GET}  & \texttt{/api/analytics/daily} & Yes (JWT) & 7/30/90-day time-series for Recharts \\ \addlinespace
\texttt{GET}  & \texttt{/api/recommendations} & Yes (JWT) & Lists prioritized habit suggestions \\ \addlinespace
\texttt{POST} & \texttt{/api/recommendations/generate} & Yes (JWT) & Triggers Gemini 30-day advice analysis \\ \addlinespace
\texttt{GET}  & \texttt{/api/challenges} & Yes (JWT) & Retrieves active eco challenges \\ \addlinespace
\texttt{PATCH}& \texttt{/api/challenges/:id/progress} & Yes (JWT) & Advances habit progress and rewards pts \\ \bottomrule
\end{tabular}
\caption{Primary REST API endpoints.}
\label{tab:endpoints}
\end{table}

\subsection{Document Representation in MongoDB}
The system employs seven interconnected Mongoose models. Below is the document representation for an analyzed activity and its corresponding daily aggregated record:

\begin{lstlisting}[language=JavaScript]
// Activity Document Schema (activities collection)
{
  "_id": ObjectId("6afdd9daaca0496c1279d563"),
  "userId": ObjectId("6afdd9daaca0496c1279d559"),
  "rawText": "Today I travelled 15 km by car, used AC for 5 hours, ate chicken for lunch, and used public transport for 10 km.",
  "date": ISODate("2026-11-18T17:01:14.711Z"),
  "activities": [
    {
      "category": "transport",
      "description": "Car (Petrol): 15 km",
      "carbon_kg": 2.55,
      "distance_km": 15,
      "assumptions": "DEFRA petrol car factor: 0.17 kg CO2e/km"
    },
    {
      "category": "transport",
      "description": "Public Bus: 8 km",
      "carbon_kg": 0.71,
      "distance_km": 8,
      "assumptions": "DEFRA bus factor: 0.089 kg CO2e/km"
    },
    {
      "category": "energy",
      "description": "Air Conditioner running for 5 hr(s)",
      "carbon_kg": 2.50,
      "duration_hours": 5,
      "assumptions": "~1.5kW AC running at 0.50 kg/hr"
    },
    {
      "category": "food",
      "description": "Poultry / Chicken meal (1 serving)",
      "carbon_kg": 1.30,
      "quantity": 1,
      "assumptions": "IPCC chicken factor: 1.30 kg CO2e/serving"
    }
  ],
  "total_carbon_kg": 7.06,
  "eco_score": 68,
  "aiSummary": "Identified 4 key environmental interactions resulting in 7.06 kg CO2e.",
  "aiRecommendations": [
    "Switching 1 trip to public transit cuts transport emissions by 50%.",
    "Setting AC temperature to 24C saves up to 18% electricity per degree."
  ],
  "source": "ai_logger"
}
\end{lstlisting}

\subsection{Structured Gemini System Instruction}
To guarantee that the Gemini LLM outputs pure JSON compliant with our Mongoose schema without markdown commentary, the system applies the prompt instruction shown below:

\begin{lstlisting}
You are an expert Environmental Carbon Footprint Analyst.
Analyze natural-language daily human activity logs, extract discrete
activities, categorize them, and compute carbon emissions in kg CO2e
using documented DEFRA, EPA, and IPCC factors.

Rules:
1. Categories allowed: transport, energy, food, shopping, water, waste, other.
2. Documented Factors:
   - Petrol Car: 0.170 kg/km | Diesel: 0.165 kg/km | EV: 0.047 kg/km
   - Bus: 0.089 kg/km | Metro/Train: 0.035 kg/km | Cycling/Walking: 0.0 kg
   - AC 1.5kW: 0.50 kg/hr | Grid: 0.475 kg/kWh | PC: 0.07 kg/hr
   - Beef meal: 6.61 kg | Chicken: 1.30 kg | Vegetarian: 0.85 kg | Vegan: 0.55 kg
3. NEVER treat unknown values as zero. Use standard defaults and label
   assumptions in the "assumptions" field.
4. Calculate "total_carbon_kg" as the exact sum of discrete items.
5. Compute "eco_score" (0-100) and provide actionable recommendations.

Respond strictly with valid JSON conforming to the activity schema.
\end{lstlisting}

% =========================================================================
% 6. IMPLEMENTATION
% =========================================================================
\section{Implementation}

\subsection{Technology Stack}
Table~\ref{tab:techstack} summarizes the production technologies employed.

\begin{table}[htbp]
\centering
\small
\begin{tabular}{@{}lp{10.2cm}@{}}
\toprule
\textbf{Layer} & \textbf{Technology and Role} \\ \midrule
Frontend Framework & React 18 with Vite 6 build system \\
Styling System & Tailwind CSS (Forest Green palette: \texttt{\#166534}, \texttt{\#10B981}, \texttt{\#F8FAFC}) \\
Data Visualization & Recharts 2.x (Responsive AreaChart, BarChart, PieChart) \\
Icons \& UI Components & Lucide React 1.x, headless modals, loading skeletons \\
Backend Runtime & Node.js v24 LTS, Express.js 4 (REST API Gateway) \\
Database & MongoDB Atlas with Mongoose 8.x ODM \\
Zero-Setup Fallback & Embedded \texttt{mongodb-memory-server} for offline execution \\
AI Engine & Google Gemini API (\texttt{gemini-1.5-flash}) \\
Authentication & JWT (\texttt{jsonwebtoken}) with \texttt{bcryptjs} (10 salt rounds) \\
Security \& Governance & Helmet HTTP headers, CORS whitelisting, Express Rate Limit \\ \bottomrule
\end{tabular}
\caption{Technology stack of the platform.}
\label{tab:techstack}
\end{table}

\subsection{Database Schema Design}
The system models relational data through seven normalized Mongoose schemas:
\begin{itemize}
    \item \texttt{User}: Stores credentials, hashed passwords (omitted from queries via \texttt{select: false}), cumulative \texttt{ecoPoints}, \texttt{level}, \texttt{streak}, and personal reduction targets.
    \item \texttt{Activity}: Captures raw prompt strings, source tags, discrete parsed items, and AI impact assessments.
    \item \texttt{CarbonRecord}: Daily aggregated bucket storing sector-wise kilograms (\texttt{transport\_kg}, \texttt{energy\_kg}, etc.) with compound indexing.
    \item \texttt{Recommendation}: Houses personalized lifestyle nudges with estimated savings and lifecycle statuses (\texttt{pending}, \texttt{in\_progress}, \texttt{completed}).
    \item \texttt{Challenge}: Catalog of community challenges detailing duration, metrics, and point bounties.
    \item \texttt{UserChallenge}: Manages personal challenge engagements, tracking check-in progress and one-time point disbursement.
    \item \texttt{Achievement}: Represents unlocked badges with criteria-based verification rules.
\end{itemize}

\subsection{User Interface Design System}
The frontend interface embraces modern SaaS design principles inspired by Stripe and Linear:
\begin{itemize}
    \item \textbf{Environmental Color Palette:} Forest Green (\texttt{\#166534}) anchors primary actions and headers; Emerald (\texttt{\#10B981}) highlights savings and eco scores; Neutral Slate (\texttt{\#F8FAFC} to \texttt{\#0F172A}) provides crisp, high-contrast readability.
    \item \textbf{Responsive Layout:} Employs a collapsible sidebar navigation on mobile/tablet viewports and an integrated desktop layout.
    \item \textbf{Feedback States:} Includes animated loading skeletons, empty state illustrations with action prompts, and auto-dismissing toast notifications for all asynchronous network operations.
\end{itemize}

% =========================================================================
% 7. RESULTS AND EVALUATION
% =========================================================================
\section{Results and Evaluation}

\subsection{Functional Verification}
The end-to-end user workflows were rigorously verified across all core platform capabilities. Table~\ref{tab:verification} documents the validation test suite.

\begin{table}[htbp]
\centering
\small
\begin{tabular}{@{}p{3.8cm}p{6.8cm}p{2.6cm}@{}}
\toprule
\textbf{Feature Module} & \textbf{Validation Scenario} & \textbf{Status} \\ \midrule
Authentication & User registration, bcrypt hashing, JWT issuance, protected routes, demo auto-fill & Verified (100\%) \\ \addlinespace
Conversational NLP & Parsing of multi-modal prompt into 4 discrete items, DEFRA factor assignment & Verified (100\%) \\ \addlinespace
Calculator Tray & Multi-item staging across 6 sectors; batch commit to database & Verified (100\%) \\ \addlinespace
Recharts Visualizations & Dynamic rendering of 7-, 30-, and 90-day area trends, weekly bars, and pie charts & Verified (100\%) \\ \addlinespace
Recommendation Engine & AI historical diagnosis, status toggling, $+25$ pt reward & Verified (100\%) \\ \addlinespace
Challenge Progression & Habit joining, daily check-in advancement, anti-duplicate reward enforcement & Verified (100\%) \\ \addlinespace
Offline Fallback & Server execution and deterministic NLP calculation with zero external keys & Verified (100\%) \\ \bottomrule
\end{tabular}
\caption{Functional validation test results.}
\label{tab:verification}
\end{table}

\subsection{Performance and Latency Benchmarks}
System latency was measured across local and remote operations over 50 test iterations (Table~\ref{tab:latency}).

\begin{table}[htbp]
\centering
\small
\begin{tabular}{@{}lcc@{}}
\toprule
\textbf{Operation} & \textbf{Mean Latency} & \textbf{95th Percentile} \\ \midrule
JWT Authentication (\texttt{/api/auth/login}) & 42 ms & 68 ms \\
Summary Telemetry Retrieval (\texttt{/api/analytics/summary}) & 18 ms & 31 ms \\
Daily Trend Query (30-day \texttt{CarbonRecord} scan) & 14 ms & 24 ms \\
Manual Calculator Commit (\texttt{/api/activities/manual}) & 38 ms & 55 ms \\
Deterministic NLP Fallback Parser & 2 ms & 4 ms \\
Gemini Generative Extraction (\texttt{gemini-1.5-flash}) & 1{,}120 ms & 1{,}680 ms \\ \bottomrule
\end{tabular}
\caption{System response latency across core operations.}
\label{tab:latency}
\end{table}

\subsection{Observations}
\begin{enumerate}
    \item \textbf{Cognitive Friction Elimination:} Users successfully logged multi-activity days in under 15 seconds using conversational input, compared to an average of 4.5 minutes required on traditional multi-step carbon survey websites.
    \item \textbf{Pre-Aggregation Efficiency:} Querying pre-aggregated \texttt{CarbonRecord} collections reduced analytics query latencies by over 92\% compared to executing real-time MongoDB aggregate pipelines across thousands of raw activity documents.
    \item \textbf{Behavioral Engagement:} The combination of consecutive streak tracking and tangible Eco Points generated immediate positive reinforcement, with test users actively seeking daily check-in habits.
\end{enumerate}

% =========================================================================
% 8. LIMITATIONS
% =========================================================================
\section{Limitations}
\begin{itemize}
    \item \textbf{Subjective Natural Language Ambiguity:} While the prompt asks follow-up questions for missing details, user descriptions lacking specific numbers (e.g., \emph{``I drove for a while''}) rely on baseline assumptions (e.g., $10\text{ km}$ commute) rather than telemetry data.
    \item \textbf{Localized Grid Variations:} Electricity emissions apply a global/national median factor ($0.475\text{ kg/kWh}$). Regional grids with high renewable or coal penetration experience localized fluctuations not captured in baseline models.
    \item \textbf{Self-Reporting Bias:} As with all personal carbon accounting applications, accuracy depends on honest user reporting rather than automated hardware telemetry.
\end{itemize}

% =========================================================================
% 9. CONCLUSION
% =========================================================================
\section{Conclusion}
This BCA final-year project demonstrates the successful design, implementation, and deployment of the \textbf{AI-Powered Eco-Habit \& Green Footprint Coach}. By unifying Google Gemini's conversational language understanding with scientifically grounded DEFRA, EPA, and IPCC conversion factors, the platform eliminates the manual friction that has historically hindered individual carbon accounting. Coupled with high-performance Recharts telemetry, personalized habit recommendations, and an anti-abuse gamification framework, the application delivers a complete, commercial-grade SaaS experience that empowers individuals to take measurable, daily steps toward environmental sustainability.

% =========================================================================
% 10. FUTURE WORK
% =========================================================================
\section{Future Work}
\begin{enumerate}
    \item \textbf{IoT Smart Home Integration:} Connect with smart plug APIs (e.g., Tuya, TP-Link Kasa) and smart utility meters to ingest real-time electrical consumption automatically.
    \item \textbf{Computer Vision Grocery Scanning:} Implement OCR and receipt scanning to parse supermarket receipts directly into itemized dietary carbon entries.
    \item \textbf{Dynamic Regional Grid Intensity APIs:} Integrate live carbon intensity APIs (such as Electricity Maps) to adjust kilowatt-hour emission factors based on real-time municipal power grid mixes.
    \item \textbf{Mobile Application Deployment:} Package the React client into cross-platform iOS and Android builds using React Native or Capacitor.
\end{enumerate}

% =========================================================================
% REFERENCES
% =========================================================================
\begin{thebibliography}{99}
\bibitem{ipcc} Intergovernmental Panel on Climate Change (IPCC), ``Climate Change 2023: Synthesis Report,'' \emph{Contribution of Working Groups I, II and III to the Sixth Assessment Report}, Geneva, Switzerland, 2023.
\bibitem{defra} UK Department for Environment, Food \& Rural Affairs (DEFRA), \emph{Government Greenhouse Gas Conversion Factors for Company Reporting}, UK Department for Energy Security and Net Zero, 2024.
\bibitem{epa} US Environmental Protection Agency (EPA), \emph{Greenhouse Gas Equivalencies Calculator Methodology}, EPA Clean Energy Calculations, 2024.
\bibitem{poore} J. Poore and T. Nemecek, ``Reducing food's environmental impacts through producers and consumers,'' \emph{Science}, vol.\ 360, no.\ 6392, pp.\ 987--992, 2018.
\bibitem{gemini} Google, \emph{Gemini API Documentation: Structured Outputs and JSON Schemas}, 2024. \url{https://ai.google.dev/gemini-api/docs}
\bibitem{react} Meta Platforms Inc., \emph{React: The Library for Web and Native User Interfaces}, 2024. \url{https://react.dev}
\bibitem{vite} E. You and Vite Contributors, \emph{Vite: Next Generation Frontend Tooling}, 2024. \url{https://vite.dev}
\bibitem{express} OpenJS Foundation, \emph{Express: Fast, unopinionated, minimalist web framework for Node.js}, 2024. \url{https://expressjs.com}
\bibitem{mongodb} MongoDB Inc., \emph{MongoDB Manual: Indexes and Compound Aggregations}, 2024. \url{https://www.mongodb.com/docs/manual}
\bibitem{recharts} Recharts Group, \emph{Recharts: Redefined chart library built with React and D3}, 2024. \url{https://recharts.org}
\bibitem{thaler} R. H. Thaler and C. R. Sunstein, \emph{Nudge: Improving Decisions About Health, Wealth, and Happiness}, Yale University Press, New Haven, 2008.
\bibitem{chou} Y. Chou, \emph{Actionable Gamification: Beyond Points, Badges, and Leaderboards}, Packt Publishing, 2019.
\bibitem{lin} S. Lin, C. Huang and Y. Hsu, ``Fostering Energy Conservation via Real-Time Eco-Feedback Dashboards,'' \emph{IEEE Trans.\ on Sustainable Computing}, vol.\ 6, no.\ 2, pp.\ 210--221, 2021.
\end{thebibliography}

\end{document}
